% Diagrama de sequência: ciclo de leitura do fluxo em nuvem.
% Firmware: Rede/transmissor_esp32cam; servidor: reconstrutor.py e pipeline_nuvem.py.
% Elementos comuns aos diagramas de sequência (uml_borda.tex e uml_nuvem.tex).
% Usado pelo texto (tcc2.tex) e pela apresentação (apresentacao_tcc2.tex).
% Requer as bibliotecas TikZ arrows.meta e calc.
%
% Posição horizontal de cada linha de vida (cm)
\def\umlXno{0}
\def\umlXrad{6.0}
\def\umlXrx{10.5}
\def\umlXsrv{15.0}
%
% \umlcabecalho{x}{texto}: caixa do participante no topo
\providecommand{\umlcabecalho}[2]{%
  \node[draw, fill=white, rounded corners=1pt, align=center, minimum width=3.2cm,
        minimum height=0.9cm, font=\small] at (#1,0) {#2};}
% \umlvida{x}{y final}: linha de vida tracejada
\providecommand{\umlvida}[2]{\draw[dashed, black!55] (#1,-0.45) -- (#1,#2);}
% \umlmsg{y}{x origem}{x destino}{rótulo}: mensagem síncrona
\providecommand{\umlmsg}[4]{%
  \draw[-{Stealth[length=2.2mm]}, thick] (#2,#1) -- node[above, font=\footnotesize,
    align=center, inner sep=2pt] {#4} (#3,#1);}
% \umlret{y}{x origem}{x destino}{rótulo}: resposta (tracejada)
\providecommand{\umlret}[4]{%
  \draw[-{Stealth[length=2.2mm, open]}, dashed] (#2,#1) -- node[above, font=\footnotesize,
    align=center, inner sep=2pt] {#4} (#3,#1);}
% \umlauto{y}{x}{rótulo}: mensagem para o próprio participante
\providecommand{\umlauto}[3]{%
  \draw[-{Stealth[length=2.2mm]}, thick] (#2,#1) -- ++(0.45,0) -- ++(0,-0.3) -- ++(-0.45,0);
  \node[anchor=west, font=\footnotesize, align=left, fill=white, inner sep=1pt]
       at (#2+0.55,#1-0.15) {#3};}
% \umlloop{y topo}{y base}{x esq}{x dir}{condição}: fragmento combinado "loop"
\providecommand{\umlloop}[5]{%
  \draw[black!70] (#3,#1) rectangle (#4,#2);
  \node[draw, fill=white, anchor=north west, font=\footnotesize\bfseries, inner sep=2pt]
       (umlL) at (#3,#1) {loop};
  \node[anchor=west, font=\footnotesize] at (umlL.east) {[#5]};}
%
\begin{tikzpicture}
  \umlcabecalho{\umlXno}{\textbf{ESP32-CAM}\\(firmware)}
  \umlcabecalho{\umlXrad}{\textbf{Radioenge RD49C}}
  \umlcabecalho{\umlXrx}{\textbf{Receptor LoRa32}}
  \umlcabecalho{\umlXsrv}{\textbf{Servidor}\\(\texttt{reconstrutor.py})}
  \foreach \x in {\umlXno,\umlXrad,\umlXrx,\umlXsrv} {\umlvida{\x}{-11.3}}

  \umlauto{-1.05}{\umlXno}{desperta; captura um quadro\\$160 \times 120$ em cinza e\\desliga a câmera}
  \umlauto{-2.55}{\umlXno}{comprime em JPEG (qualidade 90,\\2,9 a 3,4~KB) e divide em $N$\\blocos de 175 bytes}

  \umlloop{-3.5}{-8.15}{-0.4}{19.6}{para cada bloco; $N$ = 17 a 20}
  \umlmsg{-4.5}{\umlXno}{\umlXrad}{\texttt{AT+SENDB=01:}$\langle$181 bytes$\rangle$}
  \umlmsg{-5.45}{\umlXrad}{\umlXrx}{\textit{uplink} LoRaWAN\\cifrado, FPort 1}
  \umlret{-6.0}{\umlXrad}{\umlXno}{\texttt{AT\_TX\_OK}}
  \umlmsg{-6.95}{\umlXrx}{\umlXsrv}{quadro em hexadecimal\\e RSSI (USB)}
  \umlauto{-7.3}{\umlXsrv}{descriptografa (AppSKey)\\e guarda o bloco}
  \umlauto{-7.4}{\umlXno}{aguarda 7~s}

  \umlauto{-8.65}{\umlXsrv}{com todos os blocos,\\remonta e grava o JPEG}
  \umlauto{-8.65}{\umlXno}{\textit{deep sleep} até o próximo ciclo}
  \umlauto{-10.1}{\umlXsrv}{executa o OCR depois, com a\\mesma ROI calibrada à mão,\\sobre o JPEG gravado\\(\texttt{pipeline\_nuvem.py})}
\end{tikzpicture}%
